11.2 The Integral over a Set
123
defined earlier can be called (and is called) the Jordan measure of the (Jordan-
measurable) set $E$.
\subsubsection{Problems and Exercises}
\begin{enumerate}
    \item a) Show that if a set $E \subset \mathbb{R}^n$ is such that $\mu(E) = 0$, then the relation $\mu(\overline{E}) = 0$
    also holds for the closure $\overline{E}$ of the set.
    \item b) Give an example of a bounded set $E$ of Lebesgue measure zero whose closure
    $\overline{E}$ is not a set of Lebesgue measure zero.
    \item c) Determine whether assertion b) of Lemma 3 in Sect. 11.1 should be under-
    stood as asserting that the concepts of Jordan measure zero and Lebesgue measure
    zero are the same for compact sets.
    \item d) Prove that if the projection of a bounded set $E \subset \mathbb{R}^n$ onto a hyperplane $\mathbb{R}^{n-1}$
    has $(n-1)$-dimensional measure zero, then the set $E$ itself has $n$-dimensional mea-
    sure zero.
    \item e) Show that a Jordan-measurable set whose interior is empty has measure $0$.
    \item a) Is it possible for the integral of a function $f$ over a bounded set $E$, as intro-
    duced in Definition 2, to exist if $E$ is not an admissible (Jordan-measurable) set?
    \item b) Is a constant function $f : E \to \mathbb{R}$ integrable over a bounded but Jordan-
    nonmeasurable set $E$?
    \item c) Is it true that if a function $f$ is integrable over $E$, then the restriction $f|_A$ of
    this function to any subset $A \subset E$ is integrable over $A$?
    \item d) Give necessary and sufficient conditions on a function $f : E \to \mathbb{R}$ defined on
    a bounded (but not necessarily Jordan-measurable) set $E$ under which the Riemann
    integral of $f$ over $E$ exists.
    \item a) Let $E$ be a set of Lebesgue measure $0$ and $f : E \to \mathbb{R}$ a bounded continuous
    function on $E$. Is $f$ always integrable on $E$?
    \item b) Answer question a) assuming that $E$ is a set of Jordan measure zero.
    \item c) What is the value of the integral of the function $f$ in a) if it exists?
    \item The Brunn-Minkowski inequality. Given two nonempty sets $A, B \subset \mathbb{R}^n$, we form
    their (vector) sum in the sense of Minkowski $A + B := \{a + b \mid a \in A, b \in B\}$. Let
    $V(E)$ denote the content of a set $E \subset \mathbb{R}^n$.
    \item a) Verify that if $A$ and $B$ are standard $n$-dimensional intervals (parallelepipeds),
    then
    $$V^{1/n}(A + B) \ge V^{1/n}(A) + V^{1/n}(B).$$
    \item b) Now prove the preceding inequality (the Brunn–Minkowski inequality) for
    arbitrary measurable compact sets $A$ and $B$.
    \item c) Show that equality holds in the Brunn-Minkowski inequality only in the fol-
    lowing three cases: when $V(A + B) = 0$, when $A$ and $B$ are singleton (one-point)
    sets, and when $A$ and $B$ are similar convex sets.
\end{enumerate}